\documentclass[10pt,a4paper]{amsart}
\usepackage[margin=19mm]{geometry}
\usepackage{amsmath,amssymb,mathtools}
\usepackage[scr=rsfs,cal=euler]{mathalfa}
\usepackage{xfrac}
\usepackage[unicode,hidelinks,bookmarks=false]{hyperref}
\newcommand{\ie}{i.e.\ }
\newcommand{\cf}{cf.\ }
\newcommand{\resp}{respectively\ }
\newcommand{\Subsection}{\subsection}
\providecommand{\nobreakdash}{\nobreak\hbox{-}}
\setlength{\emergencystretch}{2em}
\multlinegap0pt
\usepackage{amssymb}
\usepackage{mathtools}
\usepackage[normalem]{ulem}
\newtheorem{theorem}{Theorem}
\newtheorem{corollary}{Corollary}
\newtheorem{lemma}{Lemma}
\title{Towards resurgence of Joyce structures}
\author{Iv\'an Tulli}
\date{Source: arXiv:2602.08805v1 (CC BY 4.0)}
\begin{document}
\maketitle

\noindent\emph{Source and licence.} Towards resurgence of Joyce structures. Version arXiv:2602.08805v1 (CC BY 4.0), distributed by arXiv under the Creative Commons Attribution 4.0 International licence, http://creativecommons.org/licenses/by/4.0/, as recorded in the arXiv licence metadata for this exact version and shown on the abstract page https://arxiv.org/abs/2602.08805v1. Full licence notice: \texttt{licenses/arxiv-2602.08805-CC-BY-4.0.txt}.\par\smallskip

\noindent\textit{Editorial scope.} Counted unit: the article's theorem mt2 (source0.tex lines 907--909). The article's complete original proof of it is delivered with it (source0.tex lines 911--1018, one \texttt{proof} environment of 108 lines). No mathematics is written, repaired or re-derived: the statement and the proof are the article's own lines, in the article's own notation, and no retained line is substituted or re-flowed.  Retained uncounted as the source-local context the proof needs: lines 597--599; lines 602--604; lines 896--902.  Nothing was omitted from any retained span, so the package carries no omission record.  No retained line contains a citation, so the wrapper carries no bibliography and no bibliography entry.  No retained line contains a figure environment, an image-inclusion command or drawing code, so no image is delivered and the package declares no asset dependency.  Editorial formatting only, in addition: an AMS \texttt{amsart} preamble that restates the article's own declarations and macros used by the retained lines, namely the environments \texttt{theorem}, \texttt{corollary}, \texttt{lemma} on separate counters, together with the standard packages those lines need.  The retained environments print in this document as Theorem 1, Corollary 1, Lemma 1 in order of appearance, whereas the article prints the same items with the numbering of its own full document.  The complete native source of the article as distributed in the arXiv e-print for this exact version is bundled unchanged as \texttt{sources/194276/source0.tex}.
\begin{equation}\label{maineq}
    e^{\mathcal{L}_{\dot{g}}}\mathcal{A}^{\text{st}}=\mathcal{A}\,.
\end{equation}

\begin{theorem}\label{mt1}
    Given a Joyce structure over $M$, there exists a unique $\dot{g}\in \text{Lie}(\widehat{\mathcal{G}})$ solving equation \eqref{maineq} and such that $\dot{g}|_{M}=0$.
\end{theorem}

\begin{corollary}\label{keycol} Assume that we work locally on a coordinate chart $U\subset TM$ and consider $B_r\subset B_R\subset U$ as in the previous lemma. For $l\geq 1$ the coefficients $\dot{g}_{k}$ of $\dot{g}=\sum_{k=1}^{\infty}\dot{g}_k\epsilon^k$ satisfy

\begin{equation*}
    ||\;[v_i,\dot{g}_{l+1}]\;||_{B_r} \leq \sum_{p=2}^{l+1}\frac{1}{p!}\left(\frac{2np}{R-r}\right)^p\sum_{k_1+...+k_p=l+1}\prod_{j=1}^{p}||\dot{g}_{k_j}||_{B_R}+\sum_{p=1}^{l}\frac{1}{p!}\left(\frac{2np}{R-r}\right)^p\sum_{k_1+...+k_p=l}\prod_{j=1}^{p}||\dot{g}_{k_j}||_{B_R}
\end{equation*}
for all $i=1,...,n$.
\end{corollary}

\begin{theorem}\label{mt2}
The Borel transform $\mathcal{B}[\dot{g}](\xi)$ converges in a neighbourhood of $\xi=0$ locally uniformly in $p\in TM$ if and only the same holds for $B[\dot{g}][\xi]$ when restricted to the zero section $M\subset TM$. In particular, the Borel transform of the $\dot{g}$ built in Theorem \ref{mt1} converges in a neighbourhood of $\xi=0$ locally uniformly in $p\in TM$.
\end{theorem}

\begin{proof}
    We consider first the case where $\dot{g}$ is the one built from Theorem \ref{mt1}, so that $\dot{g}|_M=0$. Furthermore, we take  $p\in M\subset TM$, a coordinate neighbourhood $U\subset TM$ of $p$ with coordinates $(Z^i,\theta^i)$ induced from coordinates $(Z^i)$ on the base $M$, and a polydisc $B_r(p)\subset U$. We will show that $\mathcal{B}[\dot{g}]$ converges in a neighbourhood of $\xi=0$, uniformly in $(Z^i,\theta^i)\in B_r$, and at the end mention how the proof can be extended to the more general cases in the statement of the theorem. \\
    
    Given a polydisc $B_R(p)\subset U$ with $r<R$, we will use the notation $a_k=|| \; \dot{g}_k \; ||_{B_r},$ $A_k=||\; \dot{g}_k\;||_{B_R} $,  and $C=\frac{2n}{R-r}$ throughout the proof.  To show what we want, it is enough to show that
    \begin{equation}\label{bt1}
        \sum_{k=1}^{\infty}\frac{a_k}{k!}\xi^{k}
    \end{equation} converges for sufficiently small $\xi$. Since $[v_i,\dot{g}_{l+1}]=\partial_{\theta^{i}}(\dot{g}_{l+1}^{k})\partial_{\theta^k}$ and $\dot{g}_{l+1}$ vanishes on the zero section,  we can estimate $\dot{g}_{l+1}(z)$ in terms of its vertical derivatives via a path integral in the vertical directions. Namely, we can  write
    \begin{equation}\label{pathintest}
        \dot{g}_{l+1}(Z,\theta)=\int_{\gamma(Z,\theta)}\partial_{\theta^i}\dot{g}(Z, \varphi)\mathrm{d}\varphi^i
    \end{equation}
    where $\gamma(Z,\theta)$ is a straight path from the $(Z,0)$ to $(Z,\theta)$, and then conclude that   \begin{equation*}
    a_{l+1}=||g_{l+1}||_{B_r}\leq r\cdot n \cdot \max\{||\;[v_i, g_{l+1}]\;||_{B_r}\}_{i=1}^{n}\,.
\end{equation*} 
Combining this inequality with the inequality from Corollary \ref{keycol} we find
\begin{equation}\label{in1}
    a_{l+1} \leq r\cdot n\left(\sum_{p=2}^{l+1}\frac{1}{p!}C^pp^p\sum_{k_1+...+k_p=l+1}\prod_{i=1}^{p}A_{k_i}+\sum_{p=1}^{l}\frac{1}{p!}C^pp^p\sum_{k_1+...+k_p=l}\prod_{i=1}^{p}A_{k_i}\right)\,.
\end{equation}
Pick $y_1>A_1$ and define  $y_{l+1}$ for $l\geq 1$ as the right hand side of the inequality \eqref{in1} with all $A_{k_i}$ replaced by the $y_{k_i}$. Note that this definition makes sense since the $y_{k_i}$ that appear have $k_i<l+1$. It then follows that $a_{k}<y_{k}$ for all $k\geq 1$. Hence, to show that \eqref{bt1} converges in a neighbourhood of $\xi=0$ it is enough to show that 
\begin{equation*}
    C(\xi):=\sum_{k=1}^{\infty}\frac{y_k}{k!}\xi^{k}=\sum_{k\geq 1}c_k\xi^{k-1}
\end{equation*}
converges in a neighbourhood of $\xi=0$. To continue we will need the following combinatorial lemma.\\

\begin{lemma}Let $p$ be a positive integer. If $k_i\geq 1$ for $i=1,...,p$ are integers satisfying $k_1+...+k_p=m$ then
\begin{equation*}
    p!\leq \frac{m!}{k_1!...k_p!}\,.
\end{equation*}
\end{lemma}
\begin{proof}
Consider first two integers $a,b\geq 1$. It is then easy to see by a direct computation that $(a+b-1)!\geq a!b!$. Successively applying this we find
\begin{equation*}
k_1!...k_p!\leq (k_1+...+k_p-p+1)!=(m-p+1)!\,.
\end{equation*}
We then obtain the desired inequality
\begin{equation*}
\frac{m!}{k_1!...k_p!}\geq \frac{m!}{(m-p+1)!}=m(m-1)...(m-p+2)\geq p(p-1)...2=p!\,.
\end{equation*}
\end{proof}
Applying the lemma for the cases $m=l+1$ and $m=l$, and using the notation $[\xi^k]C(\xi)^p$ for the coefficient of $\xi^k$ in $C(\xi)^p$. we obtain the following 
\begin{equation}\label{in2}
\begin{split}
    \frac{y_{l+1}}{(l+1)!}&=r\cdot n\left(\sum_{p=2}^{l+1}\frac{(Cp)^p}{(l+1)!p!}\sum_{k_1+...+k_p=l+1}y_{k_1}\cdot...\cdot y_{k_p}+\frac{1}{(l+1)}\sum_{p=1}^{l}\frac{(Cp)^p}{l!p!}\sum_{k_1+...+k_p=l}y_{k_1}\cdot...\cdot y_{k_p}\right)\\
    &\leq r\cdot n\left(\sum_{p=2}^{l+1}\frac{(Cp)^p}{(p!)^2}\sum_{k_1+...+k_p=l+1}\frac{y_{k_1}\cdot...\cdot y_{k_p}}{k_1!...k_p!}+\frac{1}{(l+1)}\sum_{p=1}^{l}\frac{(Cp)^p}{(p!)^2}\sum_{k_1+...+k_p=l}\frac{y_{k_1}\cdot...\cdot y_{k_p}}{k_1!...k_p!}\right)\\
    &\leq r\cdot n\left(\sum_{p=2}^{l+1}\frac{(Cp)^p}{(p!)^2}[\xi^{l+1}]C(\xi)^p+\frac{1}{(l+1)}\sum_{p=1}^{l}\frac{(Cp)^p}{(p!)^2}[\xi^l]C(\xi)^p\right)\,.\\
\end{split}
\end{equation}
We can encode \eqref{in2} for all $l\geq 1$ by writing 
\begin{equation}\label{in3}
    C(\xi)-y_1\xi \leq r\cdot n\sum_{l\geq 1}\xi^{l+1}\left(\sum_{p=2}^{l+1}\frac{(Cp)^p}{(p!)^2}[\xi^{l+1}]C(\xi)^p+\frac{1}{(l+1)}\sum_{p=1}^{l}\frac{(Cp)^p}{(p!)^2}[\xi^l]C(\xi)^p\right)\,,
\end{equation}
where $\leq$ denotes the partial order in  $\mathbb{R}[[\xi]]$ where for all  $k\geq 0$ the coefficient of $\xi^k$ on the left hand side is $\leq$ than the coefficient of $\xi^k$ on the right hands side. Now note that since $C(\xi)$ does not have a constant term, for $p>k$ we have $[\xi^k]C(\xi)^p=0$. We can then extend sums in $p$ in \eqref{in3} to infinity without changing anything.  
Swapping the order of sums we can then rewrite the above as 

\begin{equation}\label{in4}
    C(\xi)-y_1\xi \leq \sum_{p=2}^{\infty}\beta_pC(\xi)^p+\int_{0}^{\xi}\sum_{p=1}^{\infty}\beta_pC(\zeta)^p\mathrm{d}\zeta\,, \quad \beta_p=rn\frac{(Cp)^p}{(p!)^2}\,.
\end{equation}
Using that differentiation with respect to $\xi$ preserves the partial order $\leq$ on $\mathbb{R}[[\xi]]$, we find from differentiating \eqref{in4} and rearranging that 
\begin{equation*}
    C'(\xi)(1-\sum_{p\geq 2}p\beta_pC(\xi)^{p-1})\leq y_1+\sum_{p\geq 1}\beta_pC(\xi)^p\,.
\end{equation*}
To simplify the expressions we introduce the notation $\Phi_1(u):=\sum_{p\geq 2}p\beta_pu^{p-1}$ and $\Phi_2(u):=\sum_{p \geq 1}\beta_pu^p$, so that 

\begin{equation*}
    C'(\xi)(1-\Phi_1(C(\xi))\leq y_1+\Phi_2(C(\xi))\,.
\end{equation*}
Note that since $(1-\Phi_1(C(\xi)))|_{\xi=0}=1$, the  formal series $(1-\Phi_1(C(\xi)))$ is invertible in $\mathbb{R}[[\xi]].$ Furthermore $(1-\Phi_1(C(\xi)))^{-1}$ is a formal series with positive coefficients, since it is given by the formal series $\sum_{k=0}^{\infty}(\Phi_1(C(\xi)))^k$ and $\Phi_1(C(\xi))$ has positive coefficients. Since products of formal series with positive coefficients preserves the order of formal series, we obtain that 
\begin{equation}\label{in5}
    C'(\xi)\leq(1-\Phi_1(C(\xi))^{-1}(y_1+\Phi_2(C(\xi)))\,.
\end{equation}
The important thing to notice now is that the formal series $\Phi_1(u)$ and $\Phi_2(u)$ both converge. In fact they have an infinite radius of convergence. This can be seen by the Stirling approximation: as $p\to \infty$
\begin{equation}
(\beta_p)^{1/p}=(r\cdot n)^{1/p}\frac{Cp}{p!^{2/p}}\sim \frac{C e^2}{(2\pi p)^{1/p}p}\to 0\,,
\end{equation}
which shows that $\Phi_2(u)$ has infinite radius of convergence, and hence $\Phi_1(u)=\Phi_2'(u)-1$ as well. It follows that the series 
\begin{equation*}
    G(u):=(1-\Phi_1(u))^{-1}(y_1+\Phi_2(u))=\sum_{m\geq 0}g_mu^m\,.
\end{equation*}
has a positive radius of convergence, and $g_m \geq 0$ for all $m$ since it is a product of two formal series with positive coefficients.\\

Now consider the formal ODE
\begin{equation*}
    f'(\xi)=G(f(\xi)), \quad f(\xi)=0\,, \quad f(\xi)=\sum_{k\geq1}f_k\xi^k\,.
\end{equation*}
The coefficients of $f_k$ are uniquely determined by the recursive relation 
\begin{equation}\label{recursive1}
    f_1=g_0, \quad f_n=\frac{1}{n}\sum_{m=1}^{n-1}g_m\sum_{k_1+...+k_m=n-1}f_{k_1}\cdot...\cdot f_{k_m}\,.
\end{equation}
In particular, note that all coefficients are non-negative. Since $G(u)$ is analytic, the formal series $f(\xi)$ actually converges in a neighbourhood of $\xi=0$, since there is a unique analytic solution near $\xi=0$ vanishing at $\xi=0$, and the coefficients of the expansion of the analytic solution must satisfy \eqref{recursive1}.\\ 

We claim that $C(\xi)\leq f(\xi).$ First, notice that by \eqref{in5}
\begin{equation*}
    C'(0)=c_1\leq [\xi^0]G(C(\xi))=g_0, \quad f_1=f'(0)=G(0)=g_0
\end{equation*}
so $c_1\leq f_1.$ On the other hand, again by \eqref{in5}
\begin{equation*}
    2c_2=[\xi]C'(\xi)\leq [\xi]G(C(\xi))=g_1c_1\leq g_1f_1=[\xi]G(f(\xi))=[\xi]f'(\xi)=2f_2 \implies c_2\leq f_2\,.
\end{equation*}
Now assume that we have proved the result by induction up to $N$. Then by \eqref{in5}
\begin{equation*}
\begin{split}
    (N+1)c_{N+1}&=[\xi^{N}]C'(\xi)\leq [\xi^N]G(C(\xi))=\sum_{m=1}^{N}g_m\sum_{k_1+...+k_m=N}c_{k_1}\cdot...\cdot c_{k_m}\\
    &\leq \sum_{m=1}^{N}g_m\sum_{k_1+...+k_m=N}f_{k_1}\cdot...\cdot f_{k_m}=[\xi^N]G(f(\xi))=[\xi^N]f'(\xi)=(N+1)f_{N+1}
\end{split}
\end{equation*}
so $c_{N+1}\leq f_{N+1}. $ It then follows that $C(\xi)\leq f(\xi)$, and since $f(\xi)$ converges in a neighbourhood of $\xi=0$, the $C(\xi)$ must also converge in a neighbourhood of $\xi=0$. This shows that $\mathcal{B}[\dot{g}](\xi)$ converges in a neighbourhood of $\xi=0$, uniformly in $(Z^i,\theta^i)\in B_r(p)$.\\

In order to prove the more general statement of the theorem, consider $\dot{g}$ with $\dot{g}|_M=0$ as before, but assume now that $p\in TM$ is any point not necessary on the zero section. Join $p$ and $\pi(p)\in M\subset TM$ with a path on the fiber $T_{\pi(p)}M$, and cover the path by finitely many polydiscs. The main thing that changes in the proof is that when we use \eqref{pathintest} when working on a polydisc, we have an extra boundary term. We can use induction and the previous proof as the base case to show that these boundary terms contribute to convergent expressions. The same argument then works for $\dot{g}$ not necessarily vanishing on $M\subset TM$, provided we assume that the Borel transform converges when restricted to the zero section. 
\end{proof}

\end{document}
